%%%%%%%%%%%%%%%%%%%%%%%%%%%%%%%%%%%%%%%%%%%%%%%%%%%%%%%%%%%%%
% ==================== 参考文献 ====================
% ✓ 自查清单：
%   □ 是否只收录真实存在且在正文中使用 \cite{} 引用的资料？
%   □ 是否按正文首次引用顺序编号？
%   □ 作者、题名、年份、卷期、页码或网址是否完整且准确？
%   □ 是否删除未在正文引用的占位条目？
%   □ 是否未将 ChatGPT、DeepSeek、Claude、Copilot 等 AI 工具列为参考文献？
%
% 注意：thebibliography 会自动生成“参考文献”标题，本文件不要再写 \section*。

\begin{thebibliography}{99}
  \bibitem{ref-journal} 【作者】. 【文章题名】[J]. 【期刊名称】, 【年份】, 【卷】(【期】): 【起止页码】.
  \bibitem{ref-book} 【作者】. 【书名】[M]. 【出版地】: 【出版社】, 【出版年】.
  \bibitem{ref-thesis} 【作者】. 【学位论文题名】[D]. 【城市】: 【授予单位】, 【年份】.
  \bibitem{ref-conference} 【作者】. 【会议论文题名】[C]//【论文集或会议名称】. 【出版地】: 【出版者】, 【年份】: 【起止页码】.
  \bibitem{ref-online} 【作者或机构】. 【网页或数据资料题名】[EB/OL]. 【发布日期】[【引用日期】]. 【网址】.
\end{thebibliography}
